\documentclass[11pt,a4paper]{amsart}

% =============================================================================
% Clean rebuild scaffold (May 2026).
%
% Status convention: every unproven leaf is followed by a Gap remark via the
% `wip' theorem environment, carrying a `rem:wip-<name>' label so the gap is
% individually addressable from anywhere in the file.  Status is binary: a
% result is proved if and only if no `wip' remark is attached to it.  All
% bookkeeping (what is missing, where v1 prose lives, dependencies) is the
% body of that remark; nothing tracking-related lives outside the .tex.
%
% The original large draft `proof_of_rh.tex' stays in place as a quarry from
% which prose may be migrated, after restatement and refereeing, into the
% present file; see Appendix A.
% =============================================================================

\usepackage{amsmath,amssymb,amsthm,mathtools,bm,mathrsfs}
\usepackage[colorlinks=true,linkcolor=black,citecolor=black,urlcolor=black]{hyperref}
\pdfstringdefDisableCommands{%
  \def\Omega{Omega}%
  \def\Phi{Phi}%
  \def\PsiCal{Psi}%
  \def\RH{RH}%
  \def\corr{corr}%
  \def\toy{toy}%
  \def\val{val}%
  \def\aut{aut}%
  \def\jet{jet}%
  \def\far{far}%
  \def\est{est}%
  \def\rem{rem}%
  \def\sym{sym}%
  \def\curv{curv}%
  \def\tail{tail}%
  \def\mix{mix}%
  \def\core{core}%
  \def\aux{aux}%
  \def\out{out}%
  \def\pair{pair}%
  \def\lin{lin}%
  \def\new{new}%
}
\usepackage{enumitem}
\usepackage[expansion=false]{microtype}
\usepackage[margin=1in]{geometry}

\theoremstyle{plain}
\newtheorem{theorem}{Theorem}[section]
\newtheorem{proposition}[theorem]{Proposition}
\newtheorem{lemma}[theorem]{Lemma}
\newtheorem{corollary}[theorem]{Corollary}

\theoremstyle{definition}
\newtheorem{definition}[theorem]{Definition}
\newtheorem{remark}[theorem]{Remark}
\newtheorem{claim}[theorem]{Claim}
\newtheorem{hypothesis}[theorem]{Hypothesis}

\title{A Proof of the Riemann Hypothesis}

\author{John Shields}
\thanks{Preprint, clean rebuild draft.}
\date{May 2026}

\subjclass[2020]{11M26 (primary), 11M06, 11M50 (secondary)}
\keywords{Riemann hypothesis, zeta function, critical line, phase kernel,
  jet normalization, jet-Gram block, odd-moment cancellation}

% Frozen namespace, identical to proof_of_rh.tex.
\newcommand{\RH}{\mathrm{RH}}
\newcommand{\z}{\zeta}
\newcommand{\Om}{\Omega}
\newcommand{\Ph}{\Phi}
\newcommand{\Del}{\Delta}
\newcommand{\Aut}{\mathrm{aut}}
\newcommand{\sym}{\mathrm{sym}}
\newcommand{\loc}{\mathrm{loc}}
\newcommand{\sm}{\mathrm{sm}}
\newcommand{\far}{\mathrm{far}}
\newcommand{\est}{\mathrm{est}}
\newcommand{\rem}{\mathrm{rem}}
\newcommand{\toy}{\mathrm{toy}}
\newcommand{\val}{\mathrm{val}}
\newcommand{\corr}{\mathrm{corr}}
\newcommand{\main}{\mathrm{main}}
\newcommand{\err}{\mathrm{err}}
\newcommand{\jet}{\mathrm{jet}}
\newcommand{\bg}{\mathrm{bg}}
\newcommand{\curv}{\mathrm{curv}}
\newcommand{\tail}{\mathrm{tail}}
\newcommand{\mix}{\mathrm{mix}}
\newcommand{\core}{\mathrm{core}}
\newcommand{\aux}{\mathrm{aux}}
\newcommand{\out}{\mathrm{out}}
\newcommand{\pair}{\mathrm{pair}}
\newcommand{\lin}{\mathrm{lin}}
\newcommand{\new}{\mathrm{new}}
\newcommand{\off}{\mathrm{off}}
\newcommand{\odd}{\mathrm{odd}}
\newcommand{\amp}{\mathrm{amp}}
\newcommand{\actual}{\mathrm{actual}}
\newcommand{\can}{\mathrm{can}}
\newcommand{\source}{\mathrm{source}}
\newcommand{\base}{\mathrm{base}}
\newcommand{\wh}{\mathrm{wh}}
\newcommand{\phys}{\mathrm{phys}}
\newcommand{\quot}{\mathrm{quot}}
\newcommand{\Qabs}{\mathcal Q^{\mathrm{abs}}}
\newcommand{\F}{\mathrm{F}}
\newcommand{\RePart}{\operatorname{Re}}
\newcommand{\dist}{\operatorname{dist}}
\newcommand{\diag}{\operatorname{diag}}
\newcommand{\Tr}{\operatorname{Tr}}
\newcommand{\Ker}{\operatorname{Ker}}
\newcommand{\sgn}{\operatorname{sgn}}

% Skeleton-only environment: marks an unproven leaf as a remark titled
% `Gap'.  Each occurrence carries a `rem:wip-<name>' label so gaps are
% individually addressable.  Removed at journal time.
\newtheorem{wip}[theorem]{Gap}

\begin{document}

\begin{abstract}
We give a local-geometric proof of the Riemann Hypothesis.  At each candidate
off-line zero we zoom into a window of size comparable to the inverse local
height; the two-by-two jet-Gram blocks of the phase kernel at one or two
centers form a finite-dimensional comparison object, and a hypothetical
off-line zero is forced into a transverse mode that the actual zeta phase
cannot realize.  Two-point and four-point closure of this comparison gives an
unconditional zero-free strip around the critical line and hence the Riemann
Hypothesis.  An optional final section provides extra source energy from
zeta, available to the two-point and four-point closures if their standalone
forms fall short.
\end{abstract}

\maketitle

\tableofcontents

\newpage

% =============================================================================
\section{The argument in one paragraph}
\label{sec:overview}

Suppose \(\rho=\beta+i\gamma\) is a nontrivial zero of \(\zeta\) with
\(\beta\ne\tfrac12\) and \(|\gamma|\) large.  Around \(\rho\) (and its
symmetry partners \(\bar\rho,1-\rho,1-\bar\rho\)) we work in a smooth window
of size comparable to \(1/Q\), where \(Q=Q(\gamma)\) is a scale parameter
slowly growing with the height.  On this window we have a real local phase
\(\Ph\) for \(\zeta\), and the only object we ever extract from \(\Ph\) is
the phase kernel \(K_\Ph(x,y)=\sin(\Ph(x)-\Ph(y))/\pi(x-y)\).  Taking the
\(2\times 2\) blocks of \(K_\Ph\) at one or two centers and rotating into the
value/derivative basis gives finite explicit matrices: the same-point block
\(J(T)\) and the cross-block \(N_{12}(T_1,T_2)\).  Everything downstream is
finite linear algebra on these matrices.

The argument is then a comparison: a hypothetical off-line zero forces a
specific transverse mode in this finite linear-algebra picture (the
\emph{toy} mode); the actual zeta phase, after a single fixed quotient that
removes harmless value/gauge directions, cannot realize that mode.  The
clash is two-point at minimum, four-point in the symmetric quartet form, and
propagates by an odd-moment cancellation across the symmetry packet.  The
contradiction excludes off-line zeros above some height; standard low-height
verification then closes the line.

\subsection{Spine}
\label{subsec:spine}

The chain that delivers \(\RH\) is

\[
\begin{aligned}
&\text{\S\ref{sec:microscope}: local microscope}
\;\longrightarrow\;
\text{\S\ref{sec:jet-gram-limits}: jet-Gram limits}
\;\longrightarrow\;
\text{\S\ref{sec:finite-blocks}: finite-scale whitening}\\
&\;\longrightarrow\;
\text{\S\ref{sec:toy} \,/\, \S\ref{sec:zeta}: toy vs.\ \(\zeta\) comparison}
\;\longrightarrow\;
\text{\S\ref{sec:two-point} \,/\, \S\ref{sec:four-point}: pairwise / quartet rigidity}\\
&\;\longrightarrow\;
\text{\S\ref{sec:odd}: odd-moment cancellation}
\;\longrightarrow\;
\text{\S\ref{sec:zerofree}: zero-free strip}
\;\longrightarrow\;
\text{\S\ref{sec:rh}: \(\RH\)}.
\end{aligned}
\]

Section~\ref{sec:hardy} stands aside: it provides extra source energy from
\(\zeta\) that the two-point or four-point closures may consume if and only
if their standalone forms in \S\ref{sec:two-point}\,/\,\S\ref{sec:four-point}
fall short.  All cross-references to \S\ref{sec:hardy} from earlier sections
are forward references; \S\ref{sec:hardy} is closed and self-contained.

\subsection{Master conditional theorem}
\label{subsec:master}

\begin{theorem}[Master conditional]
\label{thm:master}
Assume the same-point and cross-block jet-Gram limits
(Lemmas~\ref{lem:same-point-limit} and~\ref{lem:cross-block-limit}), the
finite-scale whitening loss bound (Lemma~\ref{lem:whitening-loss}), the toy
visibility lower bound (Theorem~\ref{thm:toy-visibility}), the actual-zeta
transverse suppression bound (Hypothesis~\ref{hyp:zeta-suppression}), the
two-point or four-point closure (Theorem~\ref{thm:two-point-closure} or
Theorem~\ref{thm:four-point-closure}), the odd-moment cancellation
(Lemma~\ref{lem:odd-cancel}), and the aggregation step
(Theorem~\ref{thm:zero-free-strip}).  Then every nontrivial zero of \(\zeta\)
satisfies \(\RePart s = \tfrac12\).
\end{theorem}

\begin{wip}\label{rem:wip-master}
Statement of the conditional is firm; the inputs are the
section-level theorems below.  The proof is short --- it composes the
section-level results --- and is written out in \S\ref{sec:rh}.
\end{wip}

% =============================================================================
\section{The local microscope: phase kernels and jet coordinates}
\label{sec:microscope}

A candidate off-line zero \(\rho\) at height \(\gamma\) is invisible at any
scale much larger than \(1/\gamma\), and the analytic data of \(\zeta\)
beyond a window of width \(1/Q\) around \(\gamma\) is irrelevant to the
argument.  We therefore fix a real smooth phase \(\Ph\) for \(\zeta\) on such
a window and never use \(\zeta\) again --- only \(\Ph\), and only through one
specific kernel.  The construction below makes that reduction explicit.

\subsection{Phase and kernel}
\label{subsec:phase}

\begin{definition}[Local phase chart]
\label{def:phase-chart}
For each height \(T\) we fix a real smooth function \(\Ph_T\) on an interval
\(I_T\) of length \(\asymp 1/Q(T)\) centered at \(T\), with the following
properties:
\begin{enumerate}[label=\textup{(P\arabic*)},leftmargin=*]
\item \(\Ph_T\) is the standard real critical-line phase of \(\zeta\) on
  \(I_T\), in a chart whose denominators are bounded below by a polynomial
  in \(Q\);
\item the normalized derivative \(q:=\Ph_T'\) has a polynomial lower
  spectral bound on retained packets.
\end{enumerate}
We write \(\Ph\) and \(q\) in place of \(\Ph_T,q\) when \(T\) is fixed by
context.
\end{definition}

\begin{wip}\label{rem:wip-phase-chart}
The chart and its polynomial denominator/derivative bounds
are stated as a definition here; the existence of such a chart at every high
height, and the polynomial constants, must be proved before
Definition~\ref{def:phase-chart} is used downstream.
\end{wip}

\begin{definition}[Phase kernel]
\label{def:phase-kernel}
The phase kernel of \(\Ph\) is
\[
K_\Ph(x,y)=
\begin{cases}
\dfrac{\sin\bigl(\Ph(x)-\Ph(y)\bigr)}{\pi(x-y)}, & x\ne y,\\[1.2ex]
\dfrac{q(x)}{\pi}, & x=y,
\end{cases}
\qquad
x,y\in I_T.
\]
\end{definition}

\subsection{Point-to-jet transform}
\label{subsec:point-to-jet}

The argument never uses \(K_\Ph\) at single points; it uses \(K_\Ph\) on
small symmetric pairs \(T\pm h\) and rotates the resulting two-by-two block
out of the point basis into the value/derivative basis.  The rotation is the
fixed matrix
\[
P_h
=\frac12
\begin{pmatrix}
1 & 1\\[0.6ex]
-1/h & 1/h
\end{pmatrix}.
\]

\begin{remark}[Convention note]
\label{rem:separation-convention}
We use the symmetric pair \(T\pm h\), not \(T\pm h/2\).  All formulas in
Sections~\ref{sec:jet-gram-limits} and~\ref{sec:finite-blocks} are stated
under that convention.  Any imported lemma using the alternative convention
must be retransformed before insertion.
\end{remark}

% =============================================================================
\section{Same-point and cross-block jet-Gram limits}
\label{sec:jet-gram-limits}

The kernel \(K_\Ph\) collapses, at small \(h\), into two explicit \(2\times
2\) matrices in the value/derivative basis: a same-point block \(J(T)\)
attached to a single center, and a cross-block \(N_{12}\) attached to two
centers \(T_1,T_2\).  These are the only objects later sections manipulate;
\(\Ph\) itself never reappears past this point except through \(q\), \(q'\),
\(q''\), and \(\Del=\Ph(T_1)-\Ph(T_2)\).

\subsection{Same-point limit}
\label{subsec:same-point}

\begin{lemma}[Same-point jet-Gram limit]
\label{lem:same-point-limit}
Let \(A_h(T)\) denote the \(2\times 2\) block of \(K_\Ph\) at the points
\(T-h,T+h\).  Then
\[
P_h\,A_h(T)\,P_h^\top \;\longrightarrow\; J(T)
\qquad (h\to 0),
\]
where
\[
J(T)=\frac1\pi
\begin{pmatrix}
2q(T) & \tfrac12 q'(T)\\[1ex]
\tfrac12 q'(T) & \tfrac1{12}\bigl(q''(T)+2q(T)^3\bigr)
\end{pmatrix}.
\]
\end{lemma}

\begin{wip}\label{rem:wip-same-point-limit}
Taylor expand \(\Ph(T\pm h)\) to the order needed for
the (1,1), off-diagonal, and (2,2) entries after applying \(P_h\); verify
that singular powers of \(h\) cancel; the \(\tfrac1{12}\) coefficient and
the \(2q(T)^3\) cubic correction must be cross-checked from an independent
expansion.  v1 prose to mine: \texttt{proof\_of\_rh.tex} \S\,Local kernel and
jet normalization.
\end{wip}

\subsection{Cross-block limit}
\label{subsec:cross-block}

\begin{lemma}[Cross jet-Gram limit]
\label{lem:cross-block-limit}
For \(T_1\ne T_2\), set \(s=T_1-T_2\), \(\Del=\Ph(T_1)-\Ph(T_2)\),
\(q_j=q(T_j)\), and let \(C_h\) be the \(2\times 2\) mixed block of
\(K_\Ph\) between the pairs \(T_1\pm h\) and \(T_2\pm h\).  Then
\[
P_h\,C_h\,P_h^\top \;\longrightarrow\; \frac1\pi N_{12}(T_1,T_2)
\qquad (h\to 0),
\]
where
\[
N_{12}=
\begin{pmatrix}
\dfrac{2\sin\Del}{s}
&
\dfrac{\sin\Del - q_2 s\cos\Del}{s^2}
\\[2.4ex]
\dfrac{q_1 s\cos\Del - \sin\Del}{s^2}
&
\dfrac{(q_1+q_2)s\cos\Del + (q_1 q_2 s^2 - 2)\sin\Del}{2s^3}
\end{pmatrix}.
\]
\end{lemma}

\begin{wip}\label{rem:wip-cross-block-limit}
Four-entry Taylor expansion; sign convention on \(s\)
and \(\Del\) must be audited (the sign of \(s\) controls whether the (1,2)
or (2,1) entry carries \(q_1\) or \(q_2\)).
\end{wip}

\subsection{Jet-Gram positivity}
\label{subsec:jet-gram-positivity}

\begin{lemma}[Same-point Gram positivity]
\label{lem:same-point-positivity}
Under (P1)--(P2) of Definition~\ref{def:phase-chart}, on retained packets at
height \(T\) the block \(J(T)\) is positive definite and admits a polynomial
spectral lower bound
\[
\lambda_{\min}\bigl(J(T)\bigr) \ge Q^{-C_J}
\]
for an absolute constant \(C_J\).
\end{lemma}

\begin{wip}\label{rem:wip-same-point-positivity}
Reduce to a quadratic in \(q\), \(q'\), \(q''\);
use (P2) to bound \(q\) from below and a derivative-of-zeta-phase estimate
to bound \(q''\) from above on retained packets.  v1 prose to mine: the
relevant curvature/tail-curvature bounds.
\end{wip}

% =============================================================================
\section{Finite-scale packet blocks and whitening}
\label{sec:finite-blocks}

The \(h\to 0\) limits of \S\ref{sec:jet-gram-limits} are clean but not
directly usable: the actual proof needs the same blocks at finite separation
\(s>0\), uniformly in \(s\) and in the midpoint \(m\).  Whitening is the
change of basis that makes these finite blocks dimensionless and absorbs the
positivity of \(J(T)\) into the normalization.  The result is a single
finite-scale packet matrix \(\widehat\Om_\z(s;m)\) that is the input to
every comparison theorem from \S\ref{sec:toy} onward.

\subsection{Finite blocks}
\label{subsec:finite-blocks}

Fix a midpoint \(m\) and separation \(s>0\); set \(t_\pm=m\pm s/2\),
\(q_\pm=q(t_\pm)\), \(\Del_m(s)=\Ph(t_-)-\Ph(t_+)\).

\begin{definition}[Finite same-point and mixed blocks]
\label{def:finite-blocks}
\[
G_{m,\pm}(s)
=\frac1\pi
\begin{pmatrix}
2q(t_\pm) & \tfrac12 q'(t_\pm)\\[0.8ex]
\tfrac12 q'(t_\pm) & \tfrac1{12}\bigl(q''(t_\pm)+2q(t_\pm)^3\bigr)
\end{pmatrix},
\]
\[
N_m(s)
=\frac1\pi
\begin{pmatrix}
-\dfrac{2\sin\Del_m(s)}{s}
&
\dfrac{\sin\Del_m(s)+q_+ s\cos\Del_m(s)}{s^2}
\\[2.4ex]
-\dfrac{\sin\Del_m(s)+q_- s\cos\Del_m(s)}{s^2}
&
\dfrac{(q_-+q_+)s\cos\Del_m(s)+(2-q_- q_+ s^2)\sin\Del_m(s)}{2s^3}
\end{pmatrix}.
\]
\end{definition}

\subsection{Whitened packet block}
\label{subsec:whitening}

\begin{definition}[Whitened packet block]
\label{def:whitened-block}
Provided \(G_{m,\pm}(s)\) are positive definite, set
\[
\widehat\Om_\z(s;m)
:= G_{m,-}(s)^{-1/2}\,N_m(s)\,G_{m,+}(s)^{-1/2}.
\]
\end{definition}

\begin{lemma}[Whitening loss]
\label{lem:whitening-loss}
On retained packets there is an absolute constant \(C_W\) such that the
whitening map has operator-norm distortion at most \(Q^{C_W}\); equivalently
\[
\bigl\|G_{m,\pm}(s)^{-1/2}\bigr\|_{\op}
\;\le\; Q^{C_W}
\]
on the retained set.
\end{lemma}

\begin{wip}\label{rem:wip-whitening-loss}
Diagonal of \(G_{m,\pm}\) is bounded below by
Lemma~\ref{lem:same-point-positivity}; bound the full inverse via the
explicit \(2\times 2\) formula; uniformity in \(m\) and \(s\) is the
non-trivial part.
\end{wip}

% =============================================================================
\section{What the toy off-line model sees}
\label{sec:toy}

If \(\rho=\beta+i\gamma\) with \(\beta\ne\tfrac12\), the local geometry near
\(\rho\) is exactly that of a fixed model phase: a \emph{toy} phase
\(\Ph_\toy\) parameterized by the off-line distance \(2\beta-1\), with
explicit \(q_\toy,q_\toy',q_\toy''\).  The toy whitened block
\(\widehat\Om_\toy(s;m)\) has an explicit transverse mode --- a specific
direction in the \(2\times 2\) value/derivative space whose component does
not vanish under any of the gauge or background quotients available to
\(\zeta\).  This section states the visibility lower bound; \S\ref{sec:zeta}
states the matching suppression upper bound.

\subsection{Toy phase and toy block}
\label{subsec:toy-phase}

\begin{definition}[Toy off-line phase]
\label{def:toy-phase}
For \(\beta\in(\tfrac12,1)\) and a window scale \(Q\), the toy off-line
phase \(\Ph_\toy(t;\beta,Q)\) is the model phase of a split pair at offset
\(\pm(2\beta-1)\) on a window of size \(1/Q\), in the same coordinate and
gauge as Definition~\ref{def:phase-chart}.
\end{definition}

\begin{wip}\label{rem:wip-toy-phase}
Concrete formula for \(\Ph_\toy\) and identification of its
\(q,q',q''\) with elementary functions of \((2\beta-1)\) and the local
coordinate.  v1 prose to mine: the toy anomaly and transverse obstruction
section.
\end{wip}

\subsection{Transverse projection}
\label{subsec:transverse-projection}

\begin{definition}[Transverse projection]
\label{def:transverse-projection}
There is a finite-dimensional subspace \(\mathcal V_\val\) of the
value/derivative space spanned by the value and gauge/background
directions, and a fixed projection
\(\Pi_\trans\) onto its orthogonal complement, with the following
properties on retained packets:
\begin{enumerate}[label=\textup{(\(\Pi\)\arabic*)},leftmargin=*]
\item \(\Pi_\trans\) annihilates \(\mathcal V_\val\) exactly;
\item \(\Pi_\trans\) is bounded with polynomial loss;
\item \(\Pi_\trans\) is gauge-invariant under the residual gauge freedom
  left in \(\Ph\) after Definition~\ref{def:phase-chart}.
\end{enumerate}
The projection \(\Pi_\trans\) is fixed once and for all in this section, and
is the same operator used in the toy and the zeta packets in
\S\ref{sec:zeta}.
\end{definition}

\begin{wip}\label{rem:wip-transverse-projection}
The choice of \(\mathcal V_\val\) is the most
sensitive structural decision in the rebuild: it must precede
Theorem~\ref{thm:toy-visibility} and Hypothesis~\ref{hyp:zeta-suppression}
and not be re-chosen afterward.  v1 prose to mine: the leading value matrix
and canonical calibration section, plus the local projective rigidity
section.
\end{wip}

\subsection{Toy visibility}
\label{subsec:toy-visibility}

\begin{theorem}[Toy anomaly visibility]
\label{thm:toy-visibility}
For every off-line distance \(\beta-\tfrac12 \ne 0\) at scale \(Q\), the toy
whitened block \(\widehat\Om_\toy(s;m)\) satisfies
\[
\bigl\|\Pi_\trans\,\mathcal A_\toy\bigl(\widehat\Om_\toy(s;m)\bigr)\bigr\|
\;\ge\; Q^{-A_\toy}
\]
on a positive polynomial-weight subset of midpoints \(m\) and separations
\(s\), where \(A_\toy\) is an absolute constant and \(\mathcal A_\toy\) is
the toy anomaly functional defined explicitly in terms of
\(\widehat\Om_\toy\).
\end{theorem}

\begin{wip}\label{rem:wip-toy-visibility}
Definition of \(\mathcal A_\toy\) (a polynomial
combination of the entries of \(\widehat\Om_\toy\) selected to vanish at
\(\beta=\tfrac12\) but not for \(\beta\ne\tfrac12\)) and the lower bound
for it on retained \((m,s)\).
\end{wip}

% =============================================================================
\section{What actual \texorpdfstring{\(\zeta\)}{zeta} does instead}
\label{sec:zeta}

The actual zeta phase \(\Ph\) defines its own whitened block
\(\widehat\Om_\z(s;m)\); under the same projection \(\Pi_\trans\), it
satisfies a \emph{suppression} upper bound that is incompatible with the
toy visibility lower bound at the same exponent.  This is the load-bearing
asymmetry: the off-line scenario forces \(\Pi_\trans\widehat\Om\) to be
large; \(\zeta\) forces it to be small.

\begin{hypothesis}[Actual-zeta transverse suppression]
\label{hyp:zeta-suppression}
On retained packets, after the gauge-fixed normalization of
Definition~\ref{def:phase-chart} and the same \(\Pi_\trans\) as
Theorem~\ref{thm:toy-visibility},
\[
\bigl\|\Pi_\trans\,\mathcal A_\z\bigl(\widehat\Om_\z(s;m)\bigr)\bigr\|
\;\le\; Q^{-B_\z}
\]
for some absolute constant \(B_\z>A_\toy\), where \(\mathcal A_\z\) is the
restriction of the toy anomaly functional to the actual-zeta whitened block.
\end{hypothesis}

\begin{wip}\label{rem:wip-zeta-suppression}
This is the structural gap of the program.  The
hypothesis is consumed by Theorems~\ref{thm:two-point-closure}
and~\ref{thm:four-point-closure}.  Two routes are possible: (i) prove the
upper bound from the local jet-Gram geometry alone (the standalone route);
(ii) prove it using the source-energy reservoir of \S\ref{sec:hardy} (the
Hardy-assisted route).  v1 prose to mine: the corrected finite-\(s\)
holomorphic whitening and odd holomorphic transverse channel sections.
\end{wip}

\begin{remark}[Standalone vs.\ Hardy-assisted]
\label{rem:standalone-vs-hardy}
Hypothesis~\ref{hyp:zeta-suppression} is stated as a single object; the two
proof routes do not change its statement, only the technology of its proof.
The forward references in Theorems~\ref{thm:two-point-closure}
and~\ref{thm:four-point-closure} and in \S\ref{sec:hardy} are bookkeeping
for which route is invoked, not separate statements.
\end{remark}

% =============================================================================
\section{Two zeros at once: pairwise rigidity}
\label{sec:two-point}

The simplest comparison takes a hypothetical off-line zero \(\rho\) and a
single symmetry partner --- typically \(\bar\rho\) or \(1-\rho\) --- and
reads off the joint contribution to the cross-block \(N_m(s)\).  Under
\(\Pi_\trans\) the toy visibility bound and the actual-zeta suppression
bound apply at the same separation \(s\) and the same midpoint \(m\); they
are incompatible at the same scale exponent, and the pair cannot exist.

\begin{theorem}[Two-point closure]
\label{thm:two-point-closure}
Let \(\rho=\beta+i\gamma\) with \(\beta\ne\tfrac12\) and \(|\gamma|\) large,
and let \((T_1,T_2)\) be the height pair attached to \(\rho\) and one of its
symmetry partners.  Then under
Lemmas~\ref{lem:same-point-limit},~\ref{lem:cross-block-limit},
\ref{lem:whitening-loss}, Theorem~\ref{thm:toy-visibility}, and
Hypothesis~\ref{hyp:zeta-suppression},
\[
A_\toy < B_\z \qquad\Longrightarrow\qquad
\rho\notin\{\zeta=0\}.
\]
\end{theorem}

\begin{wip}\label{rem:wip-two-point-closure}
Compose Theorem~\ref{thm:toy-visibility} (lower) with
Hypothesis~\ref{hyp:zeta-suppression} (upper) at matched \((m,s)\); show the
exponent gap survives the polynomial whitening loss
(Lemma~\ref{lem:whitening-loss}) and the polynomial Gram lower bound
(Lemma~\ref{lem:same-point-positivity}).  This proof is short ---
arithmetic on exponents --- once its inputs are in place.
\end{wip}

\begin{remark}[Hardy assist, if needed]
\label{rem:two-point-hardy}
If Hypothesis~\ref{hyp:zeta-suppression} cannot be discharged from local
jet-Gram geometry alone for the two-point pairing, the source-energy bound
of Hypothesis~\ref{hyp:hardy-energy} in \S\ref{sec:hardy} provides the
missing room.  Theorem~\ref{thm:two-point-closure} is then conditional on
\S\ref{sec:hardy} as well as on the local inputs; the exponent inequality
\(A_\toy<B_\z\) remains the only thing to verify.
\end{remark}

% =============================================================================
\section{Four zeros in a symmetric quartet: stronger rigidity}
\label{sec:four-point}

The full symmetry packet of an off-line zero is the quartet
\[
\mathcal Z(\rho)=\{\rho,\bar\rho,1-\rho,1-\bar\rho\}.
\]
The four-center jet-Gram block carries strictly more rigidity than the
two-point version: the toy visibility lower bound holds in a refined
quotient that two-point comparisons cannot see, and the matching
suppression hypothesis is correspondingly stronger.  This section states
the four-point closure independently of \S\ref{sec:two-point}; either route
suffices for \S\ref{sec:zerofree}.

\begin{theorem}[Four-point closure]
\label{thm:four-point-closure}
Let \(\rho\) be a hypothetical off-line zero and \(\mathcal Z(\rho)\) its
symmetry quartet.  The four-center jet-Gram comparison, under the same
projection \(\Pi_\trans\) extended to the four-point quotient, yields
\(A_\toy^{(4)}<B_\z^{(4)}\) for absolute four-point exponents, and hence
\(\mathcal Z(\rho)\) cannot be a quartet of zeros.
\end{theorem}

\begin{wip}\label{rem:wip-four-point-closure}
State and prove the four-center analogue of
\(\widehat\Om_\z\), the four-point projection \(\Pi_\trans^{(4)}\), and the
four-point versions of Theorem~\ref{thm:toy-visibility} and
Hypothesis~\ref{hyp:zeta-suppression}; assemble the exponent inequality.
v1 prose to mine: the local projective rigidity of the one-pair and
multi-pair families, plus the corrected fixed-target quotient package.
\end{wip}

\begin{remark}[Hardy assist, if needed]
\label{rem:four-point-hardy}
The four-point analogue of Remark~\ref{rem:two-point-hardy} applies: if the
four-point standalone suppression cannot be proved purely locally, the
source-energy bound of Hypothesis~\ref{hyp:hardy-energy} can be invoked,
with the same effect of supplying the gap between \(A_\toy^{(4)}\) and
\(B_\z^{(4)}\).
\end{remark}

% =============================================================================
\section{Odd-moment cancellation across \texorpdfstring{\(N\)}{N} points}
\label{sec:odd}

The exponent inequality in Theorems~\ref{thm:two-point-closure}
and~\ref{thm:four-point-closure} is the load-bearing inequality, but its
inputs include a finite remainder coming from the Taylor expansion of
\(\Ph\) around midpoints.  By symmetry, odd moments of the relevant
expansion vanish under integration against the symmetric packet weight;
this is the mechanism that keeps the remainders below
\(\min(Q^{-A_\toy},Q^{-B_\z})\) and lets the comparison close.

\begin{lemma}[Odd-moment cancellation]
\label{lem:odd-cancel}
Let \(\widehat\mu_{N,Q}\) be the symmetric \(N\)-point packet measure
attached to height \(T\) and scale \(Q\).  For every \(r\ge 0\),
\[
\partial_z^{\,2r+1}\widehat\mu_{N,Q}(0)=0.
\]
The same identity holds with the packet measure replaced by its convolution
with any even bump \(\widehat\phi\): for every \(r\ge 0\),
\[
\partial_z^{\,2r+1}\bigl(\widehat\phi(\eta z)\widehat\mu_{N,Q}(z)\bigr)\big|_{z=0}=0.
\]
\end{lemma}

\begin{wip}\label{rem:wip-odd-cancel}
Pure parity argument; v1 has the full proof in the
\(N\)-point odd-moment cancellation section.  Migration is straight ---
only the symmetry packet definition needs to match
Definition~\ref{def:phase-chart}.
\end{wip}

% =============================================================================
\section{From local obstruction to a zero-free strip}
\label{sec:zerofree}

The two-point and four-point closures of \S\ref{sec:two-point} and
\S\ref{sec:four-point} exclude any single off-line zero above height
\(T_0\); to upgrade this to a zero-free strip we aggregate over heights
using the polynomial nature of every loss in
Lemmas~\ref{lem:whitening-loss},~\ref{lem:same-point-positivity},
Theorem~\ref{thm:toy-visibility}, and Hypothesis~\ref{hyp:zeta-suppression}.

\begin{theorem}[Zero-free strip]
\label{thm:zero-free-strip}
There is \(T_0>0\) such that \(\zeta(\beta+i\gamma)\ne 0\) for every
\(\beta\ne\tfrac12\) and \(|\gamma|\ge T_0\).
\end{theorem}

\begin{wip}\label{rem:wip-zero-free-strip}
Aggregate Theorem~\ref{thm:two-point-closure}
or Theorem~\ref{thm:four-point-closure} across the dyadic height ladder;
the polynomial losses combine into a single polynomial loss in \(Q(\gamma)\)
and the exponent inequality \(A_\toy<B_\z\) carries through uniformly.
\end{wip}

% =============================================================================
\section{The Riemann Hypothesis}
\label{sec:rh}

Theorem~\ref{thm:zero-free-strip} excludes high off-line zeros; standard
verification of \(\zeta\) on \(|\gamma|\le T_0\) covers the low-height
range.

\begin{theorem}[Riemann Hypothesis]
\label{thm:rh}
Every nontrivial zero of \(\zeta\) satisfies \(\RePart s = \tfrac12\).
\end{theorem}

\begin{proof}
Combine Theorem~\ref{thm:zero-free-strip} with the standard low-height
numerical verification of \(\RH\) up to height \(T_0\).
\qed
\end{proof}

% =============================================================================
\section{(Optional) Extra source energy from \texorpdfstring{\(\zeta\)}{zeta}: the Hardy reservoir}
\label{sec:hardy}

This section is a self-contained interface.  Its only role is to provide a
single source-energy lower bound that
Theorems~\ref{thm:two-point-closure} and~\ref{thm:four-point-closure} may
invoke if their standalone forms in \S\ref{sec:two-point}\,/\,\S\ref{sec:four-point}
prove insufficient.  Nothing earlier in the paper depends on this section.

\subsection{Canonical source curvature}
\label{subsec:hardy-canonical-source}

\begin{definition}[Canonical source curvature]
\label{def:canonical-source}
The canonical source curvature of \(\zeta\) on the local chart of
Definition~\ref{def:phase-chart} is
\[
q''_\can := \partial_t^3 Z_{\mathrm H}(t),
\]
where \(Z_{\mathrm H}\) is the Hardy-normalized real critical-line scalar
of \(\zeta\) in the same chart.
\end{definition}

\begin{wip}\label{rem:wip-canonical-source}
Identification of \(q''_\can\) with
\(\partial_t^3 Z_{\mathrm H}\) is convention; the chart-regularity and
zero-separation hypotheses of Definition~\ref{def:phase-chart} carry over
verbatim.  v1 prose to mine: the upstream Gate-1 arithmetic inputs section.
\end{wip}

\subsection{Source-energy reservoir}
\label{subsec:hardy-energy}

\begin{hypothesis}[Hardy source-energy bound]
\label{hyp:hardy-energy}
On retained packet intervals \(I\subset I_T\), the canonical source
curvature satisfies
\[
\int_I |q''_\can(m)|^2\,dm
\;\ge\; Q^{-A_E}\,|I|
\]
on a positive polynomial-weight subfamily, for an absolute constant
\(A_E\).
\end{hypothesis}

\begin{wip}\label{rem:wip-hardy-energy}
Hardy-type lower bound on \(\partial_t^3 Z_{\mathrm H}\)
across the retained packet family.  This is an analytic estimate on
\(\zeta\) and not on the local jet-Gram geometry; it is what makes the
section optional.  v1 prose to mine: the Hardy-curvature mass-good packet
source block.
\end{wip}

\begin{lemma}[Energy-to-suppression interface]
\label{lem:hardy-interface}
Hypothesis~\ref{hyp:hardy-energy} implies
Hypothesis~\ref{hyp:zeta-suppression} with absolute constant
\(B_\z=B_\z(A_E)\) satisfying \(B_\z>A_\toy\), via the carrier-corrected
transfer of the source energy into the transverse channel of
\(\widehat\Om_\z(s;m)\).
\end{lemma}

\begin{wip}\label{rem:wip-hardy-interface}
Carrier decomposition and packet trace step that
takes \(\int |q''_\can|^2\) on the source side into a lower bound on
\(\Pi_\trans\widehat\Om_\z\) on the jet-Gram side.  v1 prose to mine: the
canonical source-graded closure and carrier-corrected source bridge
sections.
\end{wip}

% =============================================================================
\appendix

\section{Archive map}
\label{app:archive}

The previous large draft is preserved in place at
\texttt{rh/proof\_of\_rh.tex}; it is treated as a quarry from which prose
and computations may be migrated, after restatement and refereeing, into
the present file.  Pointers to relevant v1 sections are given inline at
each \texttt{wip} mark above; the table below summarizes the main mining
correspondences.

\medskip
\noindent\begin{tabular}{@{}p{0.40\linewidth} p{0.55\linewidth}@{}}
\hline
\textbf{Rebuild section} & \textbf{v1 source (in \texttt{proof\_of\_rh.tex})} \\
\hline
\S\ref{sec:microscope} (microscope)
& Local kernel and jet normalization. \\
\S\ref{sec:jet-gram-limits} (jet-Gram limits)
& Finite-\(s\) local blocks; jet-limit local blocks. \\
\S\ref{sec:finite-blocks} (finite blocks, whitening)
& Corrected finite-\(s\) holomorphic whitening. \\
\S\ref{sec:toy} (toy visibility)
& Toy anomaly and transverse obstruction;
  leading value matrix and canonical calibration. \\
\S\ref{sec:zeta} (zeta suppression)
& Odd holomorphic transverse channel;
  corrected fixed-target quotient package. \\
\S\ref{sec:two-point}, \S\ref{sec:four-point} (two/four-point)
& Local projective rigidity of the one-pair and multi-pair families;
  finite coefficient table contracts. \\
\S\ref{sec:odd} (odd cancellation)
& \(N\)-point odd-moment cancellation. \\
\S\ref{sec:zerofree}, \S\ref{sec:rh} (strip, \(\RH\))
& Master conditional theorem, introduction. \\
\S\ref{sec:hardy} (Hardy reservoir)
& Hardy-curvature mass-good packet source block;
  carrier-corrected source bridge package;
  upstream Gate-1 arithmetic inputs. \\
\hline
\end{tabular}
\medskip

The v1 sections labeled ``Endgame architecture and remaining propagation
problem,'' ``Downstream local closure package,'' ``Unified finite-survival
architecture,'' ``Current proof state after the local package audits,''
``How the older framing maps to the new framing,'' and ``No-go and audit
ledger'' are deliberately not migrated; their content is route history and
process bookkeeping, not active proof material.

\end{document}
